\chapter{Rozwiązania projektowe}

\section{Technologia}

\begin{itemize}
    \item Implementacja została wykonana w języku C++ z wykorzystaniem standardu C++20 przy użyciu kompilatora GCC w wersji 13.3.0.
    \item Projekt budowany jest przy użyciu systemu budowania \texttt{CMake}.
    \item Testy jednostkowe zaimplementowano przy użyciu biblioteki \texttt{GoogleTest} w wersji 1.14.0.
    \item Do pomiarów wydajności wykorzystano bibliotekę \texttt{Google Benchmark} w wersji 1.9.1.
    \item Program był rozwijany i testowany w systemie operacyjnym Ubuntu 24.04.
\end{itemize}

\section{Struktura projektu}

Struktura katalogów projektu przedstawiona jest poniżej. Testy jednostkowe dla pliku \texttt{path/File.cpp}, jeśli zostały zdefiniowane,
znajdują się w katalogu \texttt{path/ut/} pod nazwą \texttt{TestFile.cpp}.

\begin{verbatim}
workspace/
|-- src/
|   |-- benchmark/              # punkt wejścia programu benchmarkującego
|   |-- common/                 # podstawowe struktury danych
|   |-- failedTests/            # szczegóły testów zakończonych niepowodzeniem
|   |-- generator/              # generatory automatów testowych
|   |-- learner/                # implementacja algorytmu ucznia
|   |-- teacher/                # implementacja nauczyciela
|   |-- utils/                  # dodatkowe narzędzia użyte w projekcie
|   |-- Tester.hpp              # klasa odpowiedzialna za uruchamianie testów
|   |-- TesterParameters.hpp    # parametry służące do konfigurowania testów
|   |-- main.cpp                # punkt wejścia programu
|-- build.sh                    # skrypt służący do budowania
|-- CMakeLists.txt              # konfiguracja systemu budowania
\end{verbatim}

\section{Struktury danych}

W tej sekcji przedstawiono kluczowe struktury danych wykorzystywane w implementacji algorytmu. Struktury te zostały zdefiniowane w module \texttt{common}.

\textbf{Alfabet wejściowy} jest reprezentowany za pomocą trzech typów wyliczeniowych: \texttt{CallSymbol}, \texttt{LocalSymbol}
oraz \texttt{ReturnSymbol}. Dzięki temu możliwe jest reprezentowanie dowolnego alfabetu
$\Sigma = \Sigma_{call} \cup \Sigma_{local} \cup \Sigma_{return}$ poprzez numerację symboli w każdym z tych zbiorów. 

Alfabet opisany jest przez krotkę $\langle |\Sigma_{call}|, |\Sigma_{local}|, |\Sigma_{return}| \rangle$
oznaczającą liczność poszczególnych zbiorów symboli.

\textbf{Zbiór symboli stosowych} jest reprezentowany za pomocą typu wyliczeniowego \texttt{StackSymbol}.
Symbol $\bot$ jest oznaczany poprzez wartość 0. Zbiór symboli stosowych musi zawierać co najmniej jeden element
reprezentujący symbol dna stosu. Jeśli $|\Sigma_{call}| > 0$ wymagane jest, żeby ten zbiór zawierał przynajmniej 2 elementy.

\textbf{Zbiór stanów automatu} również jest reprezentowany w postaci typu wyliczeniowego \texttt{State}.

\textbf{Stos} reprezentowany jest poprzez klasę \texttt{Stack} stanowiącą nakładkę na stos ze standardowej biblioteki języka C++
dostosowaną do zwracania symbolu $\bot$, gdy stos jest pusty.

\textbf{Litera kontrolna} reprezentowana jest przy użyciu tej samej struktury danych co stos automatu.
W implementacji zastosowano alias typu \texttt{ControlWord = Stack}. 

\textbf{Litera} reprezentowana jest jako typ danych \texttt{Symbol} oznaczający wariant symboli wejściowych i litery kontrolnej. 
\begin{verbatim}
std::variant<CallSymbol, ReturnSymbol, LocalSymbol, ControlWord>
\end{verbatim}

\textbf{Słowo} reprezentowane jest jako wektor liter typu \texttt{Symbol}.

\textbf{Funkcja przejścia $\delta$} jest reprezentowana jako klasa \texttt{Transition} zawierająca trzy tablice reprezentujące
przejścia dla poszczególnych typów liter z alfabetu wejściowego.
\begin{verbatim}
CoArgument callT[MaxNumOfStates][MaxNumOfLetters];
State returnT[MaxNumOfStates][MaxNumOfStackSymbols][MaxNumOfLetters];
State localT[MaxNumOfStates][MaxNumOfLetters];
\end{verbatim}

Tablice są inicjalizowane specjalnymi wartościami, które nie należą do alfabetu wejściowego,
ani zbioru symboli stosowych i reprezentują przejście do stanu \texttt{sink}.


\textbf{Deterministyczny automat visibly pushdown} reprezentowany jest jako klasa \texttt{VPA} która przechowuje zbiór stanów automatu $Q$,
stan początkowy $q_0$, zbiór stanów akceptujących $F$ i funkcję przejścia $\delta$. Alfabet wejściowy oraz zbiór symboli stosowych nie są
przechowywane jawnie w strukturze automatu, ponieważ ich reprezentacja wynika z parametrów funkcji przejścia. Wczytanie liter z poza
alfabetu nie generuje błędu, lecz prowadzi do stanu \texttt{sink}.

\textbf{String rewriting system} jest reprezentowany jako wektor reguł typu \texttt{SrsRule}, z których każda składa się z pary słów.

\section{Moduły projektu}

\subsection{Learner}
\begin{verbatim}
teacher/
|-- AutomataGenerator.hpp   # generator automatu na podstawie S i C
|-- Learner.hpp             # algorytm ucznia
|-- Selectors.hpp           # zbiór selektorów
|-- TestWords.hpp           # zbiór słów testowych
\end{verbatim}

Jest to moduł w który zawiera implementacje algorytmu ucznia. Zaczyna działanie od zainicjalizowania zbiorów $S$ i $C$,
następnie wykonuję w pętli następującą procedurę:

\begin{verbatim}
while(true)
{
    hypothesis = generator.generate();
    counterExample = oracle.equivalenceQuery(hypothesis);
    if (counterExample->empty())
        return hypothesis;
    handleCounterExample(counterExample);
}
\end{verbatim}

\textbf{Obsługa kontrprzykładu $w$} odbywa się poprzez przeiterowanie się po słowie w celu znalezienia podziału na $vaw'$.
Jeśli niezgodność zostanie wykryta na podstawie różnicy w akceptacji słowa i litera $a$ jest typu return, to zapisujemy selektor równoważny do słowa $va$,
jako wymuszony następnik automatu w konfiguracji otrzymanej po przeczytaniu słowa $v$, dla litery $a$.
Pozwala to wykluczyć to przejście, w przypadku gdyby pozostało nieokreślone na podstawie procedury generowania automatu.

\textbf{Generowanie automatu} odbywa się poprzez interfejs klasy \texttt{AutomataGenerator}. Przeplata ona procedure generowania automatu
razem z ustawianiem nieokreślonych przejść. Procedura iteruje po zbiorze selektorów konstruując funkcje przejścia, jednocześnie uzupełnia
tablice \texttt{witnesses}, w której dla każdego selektora zapisywany jest jeden świadek, który świadczy o osiągalności pewnej konfiguracji.

\begin{enumerate}
    \item Dla każdego symbolu call oraz local sprawdzane jest czy istnieje poprawny następnik w zbiorze selektorów,
    wybierając pierwszego napotkanego. Jeśli taki nie istnieje to zbiór selektorów jest rozszerzany o nowy element.

    \item Dla każdego symbolu return oraz każdego symbolu stosowego sprawdzane jest czy istnieje poprawny następnik.
    Jeśli nie istnieje to sprawdzamy, czy ta konfiguracja jest osiągalna na podstawie tablic \texttt{witnesses},
    co umożliwiłoby rozszerzenie zbioru selektorów. Jeśli nie istnieje poprawny następnik oraz świadek świadczący o osiągalności,
    to wybierany jest następnik na podstawie tablicy \texttt{forcedSelector}. Domyślnie wpisy w tej tablicy są zainicjalizowane zerami,
    co skutkuje przejściem do stanu $\epsilon$. Jednak jeśli w trakcie obsługi kontrprzykładu przejście to było uznane za błędne,
    to znajduje się w niej stan, który jest poprawnym następnikiem. 
\end{enumerate}


\subsection{Teacher}

\begin{verbatim}
teacher/
|-- Teacher.hpp           # interfejs nauczyciela
|-- Converter.hpp         # łączy automaty i konwertuje do gramatyki
|-- cfg/
    |-- Calculator.hpp    # koduje i dekoduje terminale i nieterminale
    |-- Cfg.hpp           # struktura gramatyki i sprawdzanie pustości
    |-- NonTerminal.hpp
    |-- Terminal.hpp
\end{verbatim}

Nauczyciel jako parametry wejściowe przyjmuje alfabet wejściowy, alfabet stosowy i automat docelowy. Jego zadaniem jest udostępnić interfejs,
który będzie odpowiadać na zapytania: \texttt{membershipQuery}, \texttt{stackContentQuery} i \texttt{equivalenceQuery}.
Odpowiedzi na dwa pierwsze zapytania wymagają jedynie przekazania słowa do automatu docelowego. Zapytanie o równoważność jest dominującą
częścią całego modułu. Należy podkreślić, że złożoność zapytań o równoważność w trakcie aktywnego uczenia jest znanym problemem.
Zazwyczaj jest stosowane podejście heurystyczne jak sprawdzenie losowo generowanych słów \cite{json}[5.3 Learning Algorithm],
lub bardziej złożone techniki jak np. W-method \cite{Wmethod}. Obecne podejście gwarantuje zwrócenie poprawnej odpowiedzi za każdym razem,
ale dominuje przy tym złożoność czasową i pamięciową całego programu, które wynoszą $O(?)$.

Cała procedura \texttt{equivalenceQuery} przebiega w następujących krokach:
\begin{enumerate}
    \item Oblicz $A = T \triangle H$.
    \item Przekształć $A$ na $A'$ akceptujący poprzez pusty stos.
    \item Przekształć $A'$ na gramatykę bezkontekstową $G$.
    \item Sprawdź pustość $G$.
    \item Zwróć słowo akceptowane przez $A$ na podstawie wyniku gramatyki, lub puste słowo jeśli $G$ jest pusta.
\end{enumerate}

\paragraph{Konstrukcja automatu $A = T \triangle H$}

Konstrukcja produktu automatów została zaimplementowana w metodzie \texttt{combineVPA} klasy \texttt{Converter}. Klasa ta tworzy nowy automat
symulujący jednoczesne działanie automatów $T$ oraz $H$. Zawiera on $(|Q_{T}| + 1) \cdot (|Q_{H}| + 1)$ stanów, ponieważ wszystkie przejścia
w tym automacie są zdefiniowane jawnie w przeciwieństwie do $T$ i $H$, gdzie nieokreślone przejścia były niejawnie kierowane do stanu \texttt{sink}.

\paragraph{Przekształcenie do automatu $A'$}

Przekształcenie automatu $A$ do automat akceptującego poprzez pusty stos odbywa się również w metodzie \texttt{combineVPA}.
Po wyznaczeniu stanów akceptujących możemy zdefiniować dla nich dodatkowe przejścia umożliwiające opróżnienie stosu.

\paragraph{Konstrukcja gramatyki bezkontekstowej}

Konstrukcja gramatyki bezkontekstowej $G$ została zaimplementowana w metodzie \texttt{convertVpaToCfg} klasy \texttt{Converter}.
Rozpoczyna ona od heurystycznego oszacowania liczby produkcji i nieterminali w $G$, które są podawane jako argumenty klasy \texttt{Cfg},
żeby zarezerwować pamięć na dynamicznie kontenery wymagające ciągłości pamięci.

Testy wykazały, że znacząca część dodawanych produkcji jest nieosiągalna, dlatego zastosowano heurystyczne podejście do zmniejszenia ich liczby.
Na początku dodajemy nieterminale osiągalne z nieterminala startowego razem z potrzebnymi produkcjami i umieszczamy te nieterminale w kolejce.
Następne produkcje i nieterminale się dodawne stosując algorytm BFS po osiągalnych z symbolu startowego nieterminalach.
Podejście to na losowych testach zredukowało ilość generowanych produkcji o około jedną trzecią, co również skróciło czas potrzebny na zapytanie.

\paragraph{Sprawdzanie pustości gramatyki bezkontekstowej}

Klasa \texttt{Cfg} ma na celu sprawdzić pustość poprzez wyznaczenie zbioru generujących nieterminali.
W tym celu udostępnia dwie publiczne metody: \texttt{addProduction} oraz \texttt{isEmpty}. 

Metoda \texttt{addProduction} wykorzystywana podczas konstrukcji gramatyki indeksuje nieterminale pojawiające się w produkcjach.
Dla każdej produkcji oblicza liczbę jej nieterminali po prawej stronie oraz dla każdego nieterminala zapisuję w ilu produkcjach występuje
po prawej stronie oraz zapisuję nieterminale generujące. Informacja ta pozwala skonstruować dwie tablice określające zakres produkcji zależnych
od danego nieterminala. Takie podejście umożliwia szybkie przechodzenie po wszystkich zależnych produkcjachi korzysta z ciągłości pamięci,
co w praktyce okazało się wydajne.

Sprawdzanie pustości realizowane jest w metodzie \texttt{isEmpty}. Wykonywana jest propagacja generatywności przy użyciu przeszukiwania BFS.
Jeśli wszystkie nieterminale po prawej stronie danej produkcji zostały oznaczone jako generujące, generujący staje się również jej lewy nieterminal.

Podczas działania algorytmu dla każdego generującego nieterminala zapamiętywany jest świadek jego generatywności.
Pozwala to w przypadku niepustości gramatyki odtworzyć przykładowe słowo należące do jej języka.

\paragraph{Odzyskiwanie kontrprzykładu}

Klasa \texttt{Cfg} zwraca słowo, którego prefiks rozróżnia $T$ i $H$. Ponieważ zwrócenie najkrótszego możliwego prefiksu jest
korzystne dla algorytmu uczenia, a musimy odrzucić sufiks zawierający symbole spoza alfabetu wejściowego,
przechodzi po całym słowie szukając najkrótszego prefiksu, który rozróżnia $T$ i $H$.

\subsection{Generator}

Moduł ten stanowi narzędzie służące do generowania automatów wykorzystywanych do testowania algorytmu.
Składa się z interfejsu \texttt{Generator} oraz implementacji poszczególnych generatorów, które zostały szerzej opisane w kolejnym
rozdziale~\ref{chap:generators}.

\textbf{Parametry konfiguracyjne}, które opisują generowany automat to: 

\begin{itemize}
    \item liczba stanów automatu: \texttt{numOfStates}
    \item liczba poszczególnych symboli alfabetu wejściowego: \texttt{numOfCalls}, \texttt{numOfLocals}, \texttt{numOfReturns}
    \item liczba poszczególnych symboli alfabetu stosowego: \texttt{numOfStackSymbols}
    \item gęstość przejść w automacie: \texttt{density}
    \item gęstość stanów akceptujących: \texttt{acceptingStatesDensity}
    \item liczba modułów: \texttt{numOfModules}, stosowana tylko dla automatów sterowanych symbolami call (rozdz.~\ref{chap:CDA})
\end{itemize}

\textbf{Interfejs} udostępniany przez klasę \texttt{Generator} obejmuje następujące metody:
\begin{itemize}
    \item \texttt{setConfig} umożliwiająca ustawienie parametrów konfiguracyjnych generatora.
    \item \texttt{run} generująca automat DVPA na podstawie ustawionej konfiguracji.
    \item \texttt{getSrs} zwracającą SRS dla ostatniego wygenerowanego automatu, który może być pusty jeśli generator nie wspiera generowania SRS.
\end{itemize}

\textbf{Generowanie automatu} polega na zdefiniowaniu funkcji przejścia oraz zbioru stanów akceptujących. Stan początkowy zawsze ustawiany
jest na \texttt{State\{0\}}. Ze względu na losowy charakter generatora wybór stanu początkowego nie wpływa na różnorodność generowanych automatów.
Pozostałe parametry potrzebne do konstrukcji automatu wynikają z konfiguracji generatora.

Zbiór stanów akceptujących jest wyznaczany poprzez rozpatrzenie każdego możliwego stanu i dodanie go do stanów akceptujących
z prawdopodobieństwem określonym przez parametr \texttt{acceptingStatesDensity}.

Do wygenerowania przejść w automacie dla każdego stanu, bądź konfiguracji w przypadku symbolu return losujemy czy dane przejście
będzie zdefiniowane zgodnie z parametrem \texttt{density}. Następnie spośród możliwych następników, wynikających ze specyfiki danego generatora,
losowo wybierany jest stan docelowy. W przypadku symboli typu call losowany jest również symbol stosowy z uwzględnieniem specyfiki danego generatora.

System SRS jest generowany tylko dla wybranych generatorów, gdy wynika to bezpośrednio z definicji generowanej klasy automatów DVPA.

Generator nie ustawia ziarna generatora liczb losowych, co powoduje, że kolejne uruchomienia generują różne automaty.

\subsection{Tester}

Klasa \texttt{Tester} odpowiada za przeprowadzenie zbioru testów, z których każdy polega na wygenerowaniu automatu, uruchomieniu algorytmu
aktywnego uczenia dla tego automatu oraz weryfikacji otrzymanego wyniku poprzez sprawdzenie jego wielkości oraz porównanie go z automatem
docelowym na losowo wygenerowanych słowach.

\textbf{Parametry konfiguracyjne}, które przyjmuje \texttt{Tester} to:

\begin{itemize}
    \item liczba testów do przeprowadzenia: \texttt{numOfTests}
    
    \item obiekt generatora wykorzystywany do generowania automatów testowych
    
    \item para parametrów określająca przedział, z którego losowane są parametry generatora:
    \texttt{minParameter} oraz \texttt{maxParameter} dla parametrów \texttt{numOfStates}, \texttt{numOfCalls}, \texttt{numOfLocals},
    \texttt{numOfReturns} oraz \texttt{numOfStackSymbols}
    
    \item pozostałe parametry generatora: \texttt{density}, \texttt{acceptingStatesDensity}, \texttt{numOfModules}

    \item parametr określający, czy zestaw testów powinien używać systemu SRS: \texttt{useSrs}

    \item parametry charakteryzujące maksymalną długość i liczbę słów testowych do przetestowania automatu będącego
    wynikiem działania algorytmu: \texttt{maxTestingWordLength}, \texttt{numOfRandomTestingWords}

    \item parametry służące do debugowania: \texttt{savePassedTestData}, \texttt{supervisedMode}, \texttt{supervisedTestMaxDuration}
\end{itemize}

Parametry generatora są losowo wybierane na podstawie konfiguracji. Na koniec każdego testu aktualne statystyki dotyczące przeprowadzonych
testów są drukowane na standardowe wyjście. Test kończy się niepowodzeniem, jeśli zostanie znalezione słowo, które rozróżnia wygenerowany
automat od automatu otrzymanego przez algorytm, lub jeśli otrzymany automat ma więcej stanów niż automat wygenerowany.
Jeśli test zakończy się niepowodzeniem lub parametr \texttt{savePassedTestData = true}, to szczegółowe informacje o teście
zostają zapisane w pliku \texttt{failedTests/output\_\{testNumber\}.txt}.

Użycie trybu nadzorowanego (\texttt{supervisedMode = true}) powoduje uruchomienie algorytmu uczącego w osobnym procesie.
Pozwala to kontynuować wykonywanie testów, gdy działanie programu zakończy się krytycznym błędem oraz przerwać test,
jeśli jego czas wykonania przekroczy wartość \texttt{supervisedTestMaxDuration}.

\subsection{Benchmark}
\label{chap:bench}

Moduł ten został stworzony do odpowiedniego przeprowadzenia testów na większych grupach testowych i odpowiednim zagregowaniu
wyników przy pomocy narzędzia \texttt{Google Benchmark}. TODO

\section{Podręcznik użytkownika}

W tej sekcji opisano sposób kompilacji programu oraz konfiguracji jego podstawowych parametrów.

\subsection{Uruchomienie programu}

Projekt kompilowany jest przy użyciu skryptu \texttt{build.sh}, który umożliwia:

\begin{itemize}
    \item Standardowe zbudowanie pliku wykonywalnego \texttt{build/run} przy pomocy komendy \texttt{./build.sh}
    \item Zbudowanie pliku wykonywalnego \texttt{build-debug/run} z dodatkowymi flagami debugowymi przy pomocy komendy \texttt{./build.sh -d}
    \item Zbudowanie pliku wykonywalnego służącego do benchmarkowania wyników~\ref{chap:bench} \texttt{build-bench/run} przy pomocy komendy \texttt{./build.sh -bench}
    \item Zbudowanie i uruchomienie wszystkich testów jednostkowych używając komendy \texttt{./build.sh test} (lub \texttt{./build.sh -d test})
    \item Zbudowanie i uruchomienie pojedynczego testu jednostkowego \texttt{test} używając komendy \texttt{./build.sh test -t test} (lub \texttt{./build.sh -d test -t test})
\end{itemize}

Plik wykonywalny \texttt{build/run} przyjmuje następujące argumenty:
\begin{verbatim}
./build/run numOfTest generator
\end{verbatim}
Gdzie \texttt{numOfTest} oznacza liczbę testów przekazywaną do klasy \texttt{Tester}, a \texttt{generator} określa,
który generator powinien zostać użyty i może przyjąć jedną z następujących wartości:
\textbf{random}, \textbf{cda}, \textbf{sevpa}, \textbf{mevpa}, \textbf{ecda} lub \textbf{custom}.

Podanie jako generatora \textbf{custom} ignoruje argument liczby testów i uruchamia algorytm dla jednego automatu zdefiniowanego
bezpośrednio w pliku \texttt{main.cpp} w metodzie \texttt{getCustomVpa()}. Umożliwia to zdefiniowanie własnego automatu i przetestowanie
algorytmu na jego przykładzie.

\subsection{Konfiguracja programu}

Program zawiera dwa istotne pliki konfiguracyjne:
\begin{enumerate}
    \item \texttt{utils/Constants.hpp} zawiera m.in. definicje parametrów: \texttt{MaxNumOfAutomatonStates}, \texttt{MaxNumOfLetters}
    oraz \texttt{MaxNumOfStackSymbols}, które określają maksymalny rozmiar automatów przyjmowanych przez algorytm.
    Podczas dostosowywania tych wartości należy zwrócić uwagę, że złożoność pamięciowa nauczyciela wynosi $O(?)$, a czasowa $O(?)$.

    \item \texttt{TesterParameters.hpp} zawiera definicję struktury \texttt{TesterParameters} oraz obiekty:
    \texttt{RandomTestParameters}, \texttt{CdaTestParameters}, \texttt{SeVpaTestParameters}, \texttt{MeVpaTestParameters}
    oraz \texttt{ECdaTestParameters}, które są wykorzystywane podczas odpalenia programu dla odpowiedniego generatora.
\end{enumerate}

Zmiana parametrów konfiguracyjnych wymaga ponownego skompilowania programu.